
\documentclass{article}
\usepackage{colortbl}
\usepackage{makecell}
\usepackage{multirow}
\usepackage{supertabular}

\begin{document}

\newcounter{utterance}

\centering \large Interaction Transcript for game `cladder', experiment `full\_v1.5\_default', episode 7517 with Qwen3.5{-}27B.
\vspace{24pt}

{ \footnotesize  \setcounter{utterance}{1}
\setlength{\tabcolsep}{0pt}
\begin{supertabular}{c@{$\;$}|p{.15\linewidth}@{}p{.15\linewidth}p{.15\linewidth}p{.15\linewidth}p{.15\linewidth}p{.15\linewidth}}
    \# & \multicolumn{2}{c}{Player} && \multicolumn{2}{c}{Game Master} \\
    \hline

    \theutterance \stepcounter{utterance}  
    & & & \multicolumn{4}{p{0.6\linewidth}}{
        \cellcolor[rgb]{0.9,0.9,0.9}{
            \makecell[{{p{\linewidth}}}]{
                \texttt{\tiny{[P1$\langle$GM]}}
                \texttt{You are an expert in causal inference. The following question is not a typical commonsense query, but rather a meticulously designed question created by a professor specializing in causal inference, intended to assess the students' mastery of the course content.} \\
\texttt{Imagine a self{-}contained, hypothetical world with only the following conditions, and without any unmentioned factors or causal relationships: Wibl has a direct effect on vubr and zuph. Zuph has a direct effect on uvzi. Vubr has a direct effect on uvzi. Method 1: We look at how zuph correlates with uvzi case by case according to vubr. Method 2: We look directly at how zuph correlates with uvzi in general. To understand how zuph affects uvzi, is it more correct to use the Method 1 than Method 2?} \\
\texttt{Start your answer with 'Yes' or 'No', followed by additional reasoning or evidence to support your explanation.} \\
            }
        }
    }
     \\ \\

    \theutterance \stepcounter{utterance}  
    & \multicolumn{4}{p{0.6\linewidth}}{
        \cellcolor[rgb]{0.561,0.737,0.561}{
            \makecell[{{p{\linewidth}}}]{
                \texttt{\tiny{[P1$\rangle$GM]}}
                \texttt{Yes} \\
\\ 
\texttt{In this hypothetical causal structure, **Wibl** acts as a common cause (confounder) for both **Zuph** and **Vubr**. Specifically, the causal graph contains the path \$Wibl \textbackslash{}rightarrow Zuph\$ and \$Wibl \textbackslash{}rightarrow Vubr\$. Additionally, there is a direct causal path from **Zuph** to **Uvzi** (\$Zuph \textbackslash{}rightarrow Uvzi\$) and a path from **Vubr** to **Uvzi** (\$Vubr \textbackslash{}rightarrow Uvzi\$).} \\
\\ 
\texttt{To determine the total causal effect of **Zuph** on **Uvzi**, we must consider the structure of the paths connecting them:} \\
\texttt{1.  **Direct Path**: \$Zuph \textbackslash{}rightarrow Uvzi\$ (This is the effect we want to estimate).} \\
\texttt{2.  **Backdoor Path**: \$Zuph \textbackslash{}leftarrow Wibl \textbackslash{}rightarrow Vubr \textbackslash{}rightarrow Uvzi\$.} \\
\\ 
\texttt{This backdoor path creates a non{-}causal association (confounding bias) between **Zuph** and **Uvzi** because they share a common ancestor (**Wibl**) that also influences **Vubr**, which in turn affects **Uvzi**.} \\
\\ 
\texttt{**Method 2** (looking at the general correlation) estimates the *total association*, which includes both the direct causal effect and the bias introduced by the backdoor path. Therefore, Method 2 yields a biased estimate of the causal effect.} \\
\\ 
\texttt{**Method 1** involves conditioning on **Vubr** (looking at the correlation case by case according to **Vubr**). In causal inference, to block a backdoor path, one must condition on a variable that lies on that path. Here, **Vubr** is a node on the backdoor path (\$Zuph \textbackslash{}leftarrow Wibl \textbackslash{}rightarrow Vubr \textbackslash{}rightarrow Uvzi\$). By conditioning on **Vubr**, we block the flow of non{-}causal information along this specific backdoor path. Since **Vubr** is not a collider on any path between **Zuph** and **Uvzi** (it is a mediator in the backdoor path, not a collider), conditioning on it successfully closes the backdoor without opening new spurious paths.} \\
\\ 
\texttt{Consequently, Method 1 isolates the direct causal effect of **Zuph** on **Uvzi** by removing the confounding bias, whereas Method 2 does not. Therefore, Method 1 is more correct for understanding the causal effect.} \\
            }
        }
    }
    & & \\ \\

    \theutterance \stepcounter{utterance}  
    & & & \multicolumn{2}{p{0.3\linewidth}}{
        \cellcolor[rgb]{0.9,0.9,0.9}{
            \makecell[{{p{\linewidth}}}]{
                \texttt{\tiny{[GM$|$GM]}}
                \texttt{Yes} \\
            }
        }
    }
    & & \\ \\

    \theutterance \stepcounter{utterance}  
    & & & \multicolumn{2}{p{0.3\linewidth}}{
        \cellcolor[rgb]{0.9,0.9,0.9}{
            \makecell[{{p{\linewidth}}}]{
                \texttt{\tiny{[GM$|$GM]}}
                \texttt{yes} \\
            }
        }
    }
    & & \\ \\

    \theutterance \stepcounter{utterance}  
    & & & \multicolumn{2}{p{0.3\linewidth}}{
        \cellcolor[rgb]{0.9,0.9,0.9}{
            \makecell[{{p{\linewidth}}}]{
                \texttt{\tiny{[GM$|$GM]}}
                \texttt{game\_result = WIN} \\
            }
        }
    }
    & & \\ \\

\end{supertabular}
}

\end{document}
